% FortySecondsCV LaTeX template
% Copyright © 2019-2020 René Wirnata <rene.wirnata@pandascience.net>
% Licensed under the 3-Clause BSD License. See LICENSE file for details.
%
% Please visit https://github.com/PandaScience/FortySecondsCV for the most
% recent version! For bugs or feature requests, please open a new issue on
% github.
%
% Contributors
% ------------
% * ifokkema
% * Bertbk
% * Hespe
%
% Attributions
% ------------
% * fortysecondscv is based on the twentysecondcv class by Carmine Spagnuolo
%   (cspagnuolo@unisa.it), released under the MIT license and available under
%   https://github.com/spagnuolocarmine/TwentySecondsCurriculumVitae-LaTex
% * further attributions are indicated immediately before corresponding code


%-------------------------------------------------------------------------------
%                             ADDITIONAL PACKAGES
%-------------------------------------------------------------------------------
\documentclass[
	a4paper,
	% showframes,
	% vline=2.2em,
	maincolor=cvgreen,
	% sidecolor=gray!50,
	% sectioncolor=red,
	% subsectioncolor=orange,
	% itemtextcolor=black!80,
	% sidebarwidth=0.4\paperwidth,
	% topbottommargin=0.03\paperheight,
	% leftrightmargin=20pt,
	% profilepicsize=4.5cm,
	% profilepicborderwidth=3.5pt,
	% profilepicstyle=profilecircle,
	% profilepiczoom=1.0,
	% profilepicxshift=0mm,
	% profilepicyshift=0mm,
	% profilepicrounding=1.0cm,
]{fortysecondscv}

\usepackage{ae,lmodern}
\usepackage[french]{babel}  % package pour langue française
\usepackage{graphicx} % insertion d'images et manipulation de la position
\usepackage{tikz} % création de schémas directement dans les slides

% improve word spacing and hyphenation
\usepackage{microtype}
\usepackage{ragged2e} % justification du texte dans les itemizes

% uncomment in case you don't want any hyphenation
% \usepackage[none]{hyphenat}

% take care of proper font encoding
\ifxetexorluatex
	\usepackage{fontspec}
	\defaultfontfeatures{Ligatures=TeX}
%	\newfontfamily\headingfont[Path = fonts/]{segoeuib.ttf} % local font
\else
    \usepackage[utf8]{inputenc} % compilation des mots accentués
    \usepackage[T1]{fontenc}      % césure correcte des mots accentués 
%	\usepackage[sfdefault]{noto} % use noto google font
\fi

% enable mathematical syntax for some symbols like \varnothing
\usepackage{amssymb} % caractères mathématiques

% bubble diagram configuration
\usepackage{smartdiagram}
\smartdiagramset{
	% default font size is \large, so adjust to harmonize with sidebar layout
	bubble center node font = \footnotesize,
	bubble node font = \footnotesize,
	% default: 4cm/2.5cm; make minimum diameter relative to sidebar size
	bubble center node size = 0.4\sidebartextwidth,
	bubble node size = 0.25\sidebartextwidth,
	distance center/other bubbles = 1.5em,
	% set center bubble color
	bubble center node color = maincolor!70,
	% define the list of colors usable in the diagram
	set color list = {maincolor!10, maincolor!40,
	maincolor!20, maincolor!60, maincolor!35},
	% sets the opacity at which the bubbles are shown
	bubble fill opacity = 0.8,
}


%-------------------------------------------------------------------------------
%                            PERSONAL INFORMATION
%-------------------------------------------------------------------------------
%% mandatory information
% your name
\cvname{William PENSEC}
% job title/career
\cvjobtitle{Postdoctorant en Informatique}

%% optional information
% profile picture
% \cvprofilepic{pics/profile.png}

% NOTE: ordering in sidebar will mimic the following order
% date of birth
\cvbirthday{26 octobre 1996}
% short address/location, use \newline if more than 1 line is required
\cvaddress{42000 Saint-Etienne}
% phone number
% \cvphone{+33 06.51.36.06.63}
% personal website
\cvsite{www.pensec.fr/}
% email address
\cvmail{william@pensec.fr}
% driver license
\cvcustomdata{\faCar}{Permis B - Véhiculé}
% any other custom entry
\cvcustomdata{\faMap}{Français}

%-------------------------------------------------------------------------------
%                              SIDEBAR 1st PAGE
%-------------------------------------------------------------------------------
% add more profile sections to sidebar on first page
\addtofrontsidebar{
	% include gosquare national flags from https://github.com/gosquared/flags;
	% naming according to ISO 3166-1 alpha-2 country codes
	\graphicspath{{pics/flags/}}

	% social network accounts incl. proper hyperlinks
	\profilesection{Réseaux}
		\begin{icontable}{2.5em}{1em}
			\social{\faLinkedin}
				{https://www.linkedin.com/in/william-pensec-7464a217b/}
				{LinkedIn}
            \social{\aiGoogleScholar}
                {https://scholar.google.com/citations?user=AJE3er8AAAAJ}
                {Google Scholar}
			\social{\faGithub}
				{https://github.com/WilliamPsc/}
				{WilliamPsc}
            % \social{\faUniversity}
            %     {https://labsticc.fr/fr/annuaire/pensec-william}
            %     {Lab-STICC}
		\end{icontable}

	\profilesection{Langues}
    	\pointskill{\flag{FR.png}}{Français}{5}
    	\pointskill{\flag{GB.png}}{Anglais}{4}
		\pointskill{\flag{IT.png}}{Italien}{2}
		% \pointskill{\flag{ES.png}}{Espagnol}{2}

	\profilesection{Informatique}
		\pointskill{\faCode}{Développement logiciel}{3}[4]
		\skill{}{C, Python, C++, Java, JavaScript, Android, Bash, PHP, Ada}%, Visual Basic}
		\pointskill{\faCogs}{Développement matériel}{3}[4]
		\skill{}{VHDL, System Verilog, Suite Vivado HLS, Questasim, Assembleur, RISC-V}
		\pointskill{\faDatabase}{Gestion des données}{4}[4]
		\skill{}{MySQL, JSON, XML}
		\pointskill{\faLinux}{Systèmes d'exploitation}{4}[4]
		\skill{}{Linux, Windows}
		\pointskill{\faTerminal}{Microcontrôleurs}{3}[4]
		\skill{}{Arduino, Raspberry Pi, ZedBoard}
		\pointskill{\faTerminal}{Divers}{4}[4]
		\skill{}{Latex, TikZ, Chip Whisperer}

	\profilesection{Centres d'intérêt}
		\pointskill{\faHome}{Sport}{1}[1]
			\skill[1.8em]{}{Natation, vélo, activités nautiques, randonnées}
		\pointskill{\faHome}{Loisirs}{1}[1]
			\skill[1.8em]{}{Voyages, jeux vidéos, films, séries, musique, lecture}
}


% %-------------------------------------------------------------------------------
% %                              SIDEBAR 2nd PAGE
% %-------------------------------------------------------------------------------
% \addtobacksidebar{
% 	\profilesection{About Me}
% 	\aboutme{
% 		The giant panda is a terrestrial animal and primarily spends its life
% 		roaming and feeding in the bamboo forests of the Qinling Mountains and in
% 		the hilly province of Sichuan.
% 	}

% 	\profilesection{Diagrams}
% 	\begin{sidebarminipage}
% 		\chartlabel{Bubble}
% 		\chartlabel{Diagrams}
% 		\chartlabel{with}
% 		\chartlabel{proper}
% 		\chartlabel{overflow}
% 		\chartlabel{protection}
% 		\chartlabel{for}
% 		\chartlabel{labels}
% 	\end{sidebarminipage}

% 	\begin{figure}\centering
% 		\smartdiagram[bubble diagram]{
% 			\textcolor{white}{\textbf{Being a}} \\
% 			\textcolor{white}{\textbf{Panda}}, % center bubble
% 			\textcolor{black!90}{Eating},
% 			\textcolor{black!90}{Sleeping},
% 			\textcolor{black!90}{Rolling},
% 			\textcolor{black!90}{Playing},
% 			\textcolor{black!90}{Chilling}
% 		}
% 	\end{figure}

% 	\chartlabel{Wheel Chart}

% 	\wheelchart{3.7em}{2em}{%
% 	20/3em/maincolor!50/Chill,
% 	15/3em/maincolor!15/Play,
% 	30/4em/maincolor!40/Sleep,
% 	20/3em/maincolor!20/Eat
% 	}

% 	\profilesection{Barskills}
% 	\barskill{\faSkyatlas}{Wearing asian rice hats}{60}
% 	\barskill{\faImage}{Playing Chess}{30}
% 	\barskill{\faMusic}{Playing the bamboo flute}{50}

% 	\profilesection{Memberships}
% 	\begin{memberships}
% 		\membership[4em]{pics/logo.png}{PandaScience.net}
% 		\membership[4em]{pics/logo.png}{Some longer text spanning over more than
% 			only one line}
% 	\end{memberships}
% }


%-------------------------------------------------------------------------------
%                         TABLE ENTRIES RIGHT COLUMN
%-------------------------------------------------------------------------------
\begin{document}

\makefrontsidebar

\cvsection{Profil}
\aboutme{
    Actuellement, postdoctorant en Informatique, au Laboratoire Hubert Curien à Saint-Etienne, spécialisé en sécurité matérielle sur RISC-V. Je cherche un poste de Maitre de Conférence pour poursuivre ma carrière professionnelle et approfondir mes recherches dans l'objectif de m'ouvrir à de nouvelles applications de la sécurité matérielle.
}

\cvsection{Expériences professionnelles}
\begin{cvtable}[2]
        \cvitem{2024 (1 an)}{Postdoctorat en Informatique}{Laboratoire Hubert Curien - Saint-Etienne}{Évaluation de la sécurité des implémentations FPGA de réseaux neuronaux sur RISC-V}
        \cvitem{2023 (5 mois)}{Mobilité Internationale}{ALaRI - Lugano, Suisse}{\'Etude et développement de contremesures pour la protection d'un système contre les attaques par injections de fautes (SystemVerilog).}
        \cvitem{2021 (5 mois)}{Stage de fin d'études (M2)}{Lab-STICC}{Développement d'une coopération de drones dans un système hétérogène (C++, Java, TCP/IP, CNN embarqué).}
        \cvitem{2020 (2 mois)}{Stage en Recherche (M1)}{Lab-STICC}{Simulation de drones avec capteurs (GPS et accéléromètre) dans un environnement ouvert (C++, CARES).}
\end{cvtable}


\cvsection{Formation}
\begin{cvtable}[2]
        \cvitem{2021 -- 2024}{Thèse de Doctorat (Lab-STICC)}{Université Bretagne Sud - Lorient}{Extension de la Protection des Processeurs Contre les Menaces Physiques et Logicielles par la Sécurisation du Mécanisme DIFT Contre les Attaques par Injections de Fautes - Label Européen}
        \cvitem{2019 -- 2021}{Master Logiciel pour les Systèmes Embarqués}{Université de Bretagne Occidentale - Brest}{}
        \cvitem{2015 -- 2019}{Licence en Informatique : Fondements et Applications}{Université de Bretagne Occidentale - Brest}{}
        \cvitem{2014 -- 2015}{PACES}{Université de Bretagne Occidentale - Brest}{}
        \cvitem{2014}{Baccalauréat S-SVT, spécialité ISN}{Lycée de Cornouaille - Quimper}{}
\end{cvtable}

\cvsection{Encadrements}
\begin{cvtable}[2]
        \cvitem{2022}{Encadrement d'un étudiant de M2 en projet}{Université Bretagne Sud}{Implémentation d'un processeur RISC-V sur une cible FPGA}
        \cvitem{2022}{Encadrement d'un stagiaire de M1}{Université Bretagne Sud}{Attaque physique sur cible FPGA}
        \cvitem{2021 -- 2024}{Cours effectués}{Université Bretagne Sud}{TD à hauteur de 136~heures entre la L1 (IUT) et M2.}
        % \cvitem{2020 -- 2021}{Projet de recherche (M2)}{Université de Bretagne Occidentale}{Création d'un simulateur de drone en 3D à partir de OMNeT++ et INET (C++, XML)}
        % \cvitem{2020}{Projet Robotique (M1)}{Université de Bretagne Occidentale}{Création d'un environnement pour le suivi d'un robot à distance (NodeJS, TCP/IP, OpenStreetMap)}
\end{cvtable}

\cvsection{Publications}
\begin{cvtable}
	\cvpubitem{Scripting the Unpredictable: Automate Fault Injection in CABA Simulation for Vulnerability Assessment}{W. PENSEC, V. LAPÔTRE and G. GOGNIAT}{DSD}{2024}
	\cvpubitem{Defending the Citadel: Fault Injection Attacks against Dynamic Information Flow Tracking and Related Countermeasures}{W. PENSEC, F. REGAZZONI, V. LAPÔTRE and G. GOGNIAT}{ISVLSI}{2024}
	\cvpubitem{\href{https://doi.org/10.1145/3628356.3630116}{Another Break in the Wall: Harnessing Fault Injection Attacks to Penetrate Software Fortresses}}{W. PENSEC, V. LAPÔTRE and G. GOGNIAT}{Sensors S\&P}{2023}
	\cvpubitem{\href{https://ieeexplore.ieee.org/document/9836057}{Smart Anomaly Detection and Monitoring of Industry 4.0 by Drones}}{W. PENSEC, D. ESPES, C. DEZAN}{ICUAS}{2022}
\end{cvtable}


% \newpage
% \makebacksidebar
% % \newgeometry{
% % 	top=\topbottommargin,
% % 	bottom=\topbottommargin,
% % 	right=\leftrightmargin,
% % 	left=\leftrightmargin
% % }

% \cvsection{section}
% \cvsubsection{Subsection}
% \begin{cvtable}
% 	\cvitem{<dates>}{<cv-item title>}{<location>}{<optional: description>}
% \end{cvtable}

% \cvsection{cvitem}
% \cvsubsection{Multi-line with longer description}
% \begin{cvtable}
% 	\cvitem{date}{Description}{location}{Some longer and more detailed
% 		description, that takes two lines of space instead of only one.}
% 	\cvitem{date}{Description}{location}{Some longer and more detailed
% 		description, that takes two lines of space instead of only one.}
% 	\cvitem{date}{Description}{location}{Some longer and more detailed
% 		description, that takes two lines of space instead of only one.}
% \end{cvtable}

% \cvsubsection{One-line without description}
% \begin{cvtable}
% 	\cvitem{Award}{One-line description}{Sponsor}{}
% 	\cvitem{Award}{One-line description}{Sponsor}{}
% 	\cvitem{Award}{One-line description}{Sponsor}{}
% \end{cvtable}

% \cvsection{cvitemshort}
% \cvsubsection{One-line}
% \begin{cvtable}
% 	\cvitemshort{Key}{Some further description}
% 	\cvitemshort{Key}{Some further description}
% 	\cvitemshort{Key}{Some further description}
% \end{cvtable}

% \cvsubsection{Multi-line with longer description}
% \begin{cvtable}
% 	\cvitemshort{Key}{Some further description. Can fill even more than
% 		only one single line while still keeping the correct indendation level.}
% 	\cvitemshort{Key}{Some further description. Can fill even more than
% 		only one single line while still keeping the correct indendation level.}
% 	\cvitemshort{Key}{Some further description. Can fill even more than
% 		only one single line while still keeping the correct indendation level.}
% \end{cvtable}

% \cvsection{cvpubitem}
% \begin{cvtable}
% 	\cvpubitem{Publication title}{Authors}{Journal}{Year}
% 	\cvpubitem{Publication title}{Authors}{Journal}{Year}
% 	\cvpubitem{Publication title that is spanning over multiple lines and still
% 		does not look too bad}{Authors}{Journal}{Year}
% \end{cvtable}

\cvsignature
\clearpage

\end{document}
